% !TEX program = xelatex
% --- UNIVERSAL PREAMBLE BLOCK ---
\documentclass[11pt, a4paper]{article}

% --- PACKAGES FOR MODERN WHITEPAPER STYLE ---
% Load xcolor with table option FIRST to ensure rowcolor works
% This MUST be the very first package loaded to avoid conflicts
\usepackage[table]{xcolor} 

\usepackage[a4paper, top=2.5cm, bottom=2.5cm, left=2.5cm, right=2.5cm]{geometry}
\usepackage{fontspec}

% Language Setting - Fixed for Traditional Chinese (Taiwan)
\usepackage{xeCJK}
\usepackage[provide=*]{babel}
\babelprovide[import, main]{chinese-traditional}
\babelprovide[import]{english}

% 設定英文字體
\setmainfont{Arial}[
  BoldFont=Arial Bold,
  ItalicFont=Arial Italic
]
\setsansfont{Arial}
\setmonofont{Menlo}

% 設定中文字體（使用 xeCJK）
\setCJKmainfont{PingFang TC}
\setCJKsansfont{PingFang TC}
\setCJKmonofont{PingFang TC}

% --- OTHER PACKAGES ---
% --- OTHER PACKAGES ---
% 在載入 tcolorbox 前完全禁用標籤輸出，以修復 TeX Live 2025 相容性問題
\makeatletter
% 保存原始命令
\let\orig@Hy@raisedlink\Hy@raisedlink
\let\orig@hyper@anchorstart\hyper@anchorstart
\let\orig@hyper@anchorend\hyper@anchorend

% 暫時禁用所有可能產生輸出的命令
\def\Hy@raisedlink#1{}
\def\hyper@anchorstart#1{}
\def\hyper@anchorend{}

% 重定義 tcolorbox 的標籤命令為完全不產生輸出 (Fix Undefined Control Sequence)
\def\NewStructureName#1{\@bsphack\@esphack}
\def\AssignStructureRole#1#2{\@bsphack\@esphack}
\def\NewTaggingSocket#1#2{\@bsphack\@esphack}
\def\NewTaggingSocketPlug#1#2{\@bsphack\@esphack}
\def\AssignTaggingSocketPlug#1#2{\@bsphack\@esphack}
\def\UseStructureName#1{}
\def\tagstructbegin#1{}
\def\tagstructend{}
\def\tagpdfparaOff#1{}

\makeatother

\usepackage{tcolorbox}

% 恢復原始命令
\makeatletter
\let\Hy@raisedlink\orig@Hy@raisedlink
\let\hyper@anchorstart\orig@hyper@anchorstart
\let\hyper@anchorend\orig@hyper@anchorend
\makeatother


\usepackage{atbegshi} % 用於控制頁面輸出
\usepackage{tikz}
\usepackage{listings}
\usepackage{booktabs}
\usepackage{array}
\usepackage{enumitem}
\usepackage{amssymb}
\usepackage{hyperref}
\usepackage{float}
\usepackage{longtable}

% 設定 hyperref 不顯示連結邊框
\hypersetup{
    colorlinks=false,        % 不使用彩色文字
    hidelinks,               % 完全隱藏連結（無邊框、無顏色）
    pdfborder={0 0 0},      % 設定邊框寬度為 0
}

% --- 目錄中文化設定 ---
\usepackage{titletoc}

% 設定目錄標題
\renewcommand{\contentsname}{目錄}
\renewcommand{\listfigurename}{圖目錄}
\renewcommand{\listtablename}{表目錄}

% 設定目錄格式
\titlecontents{section}
  [0em]
  {\addvspace{1em}\bfseries\sffamily\color{NavyBlue}}
  {\contentslabel{2em}}
  {}
  {\titlerule*[0.5pc]{.}\contentspage}

\titlecontents{subsection}
  [2em]
  {\addvspace{0.5em}\sffamily}
  {\contentslabel{2.5em}}
  {}
  {\titlerule*[0.5pc]{.}\contentspage}

% --- COLOR PALETTE ---
\definecolor{NavyBlue}{RGB}{0, 31, 63}
\definecolor{ElectricCyan}{RGB}{57, 204, 204}
\definecolor{CoralRed}{RGB}{255, 65, 54}
\definecolor{LightGray}{RGB}{245, 245, 245}
\definecolor{CodeBg}{RGB}{40, 44, 52}
\definecolor{SuccessGreen}{RGB}{46, 204, 64}

% --- TCOLORBOX LIBRARIES ---
\tcbuselibrary{skins, breakable, raster} 

% --- STYLING DEFINITIONS ---

% 1. Code Block Style (Dark Mode IDE)
\lstset{
    language=Java, 
    basicstyle=\ttfamily\small\color{white},
    keywordstyle=\color{ElectricCyan}\bfseries,
    commentstyle=\color{gray},
    stringstyle=\color{CoralRed},
    numberstyle=\tiny\color{gray},
    numbers=left,
    stepnumber=1,
    frame=none,
    breaklines=true,
    backgroundcolor=\color{CodeBg},
    tabsize=2,
    captionpos=b,
    morekeywords={async, await, const, let, var, function, return, import, from, export, default, describe, it, expect}
}

\newtcolorbox{codebox}[2][]{
    enhanced,
    colback=CodeBg,
    colframe=NavyBlue,
    arc=4pt,
    boxrule=0.5pt,
    title={#2},
    coltitle=white,
    fonttitle=\bfseries\sffamily,
    attach boxed title to top left={yshift=-2mm, xshift=2mm},
    boxed title style={colback=NavyBlue},
    #1
}

% 2. Callout Block (For Definitions)
\newtcolorbox{calloutbox}[2][]{
    enhanced,
    colback=ElectricCyan!10!white,
    colframe=ElectricCyan,
    leftrule=4pt,
    rightrule=0pt,
    toprule=0pt,
    bottomrule=0pt,
    arc=0pt,
    fonttitle=\bfseries\sffamily,
    coltitle=NavyBlue,
    title={#2},
    #1
}

% 3. KPI Card (For Coverage Stats)
\newtcolorbox{kpicard}[2][]{
    enhanced,
    colback=white,
    colframe=LightGray,
    coltitle=NavyBlue,
    fonttitle=\bfseries\large\centering,
    title={#2},
    halign=center,
    valign=center,
    drop shadow,
    #1
}

% 4. Section Styling
\usepackage{titlesec}
\titleformat{\section}
  {\color{NavyBlue}\normalfont\Large\bfseries\sffamily}
  {\thesection}{1em}{}[{\titlerule[1.5pt]}]
\titleformat{\subsection}
  {\color{NavyBlue}\normalfont\large\bfseries\sffamily}
  {\thesubsection}{1em}{}

% --- DOCUMENT BEGINS ---
\begin{document}

% --- COVER PAGE ---
\begin{titlepage}
  \begin{tikzpicture}[remember picture, overlay]
    % Background
    \fill[NavyBlue] (current page.south west) rectangle (current page.north east);
    % Decorative Circle
    \draw[ElectricCyan, opacity=0.3, line width=2pt] (current page.center) circle (10cm);
    \draw[CoralRed, opacity=0.3, line width=1pt] (current page.center) circle (8cm);
  \end{tikzpicture}

  \centering
  \vspace*{4cm}

  {\Huge \bfseries \sffamily \color{white} 軟體測試專題報告 \par}
  \vspace{1cm}
  {\Large \sffamily \color{ElectricCyan} FCU Sigmaboy 物品交換平台 \par}
  \vspace{0.5cm}
  {\large \sffamily \color{white} 前端應用程式軟體測試技術白皮書 \par}

  \vfill

  \begin{tcolorbox}[colback=white!10!NavyBlue, colframe=ElectricCyan, width=0.8\textwidth, arc=2mm]
    \centering \color{white}
    \textbf{專題組員：} FCU Sigmaboy Team \\
    \textbf{發布日期：} 2026-01-02 \\
    \textbf{測試核心：} API 層單元測試 (Vitest)
  \end{tcolorbox}

  \vspace{2cm}
\end{titlepage}

% --- TABLE OF CONTENTS ---
\newpage
\begin{center}
  {\Huge\bfseries\sffamily\color{NavyBlue} 目錄}
  \vspace{1cm}
\end{center}
\tableofcontents
\newpage

% --- CONTENT ---

\section{專題描述}

\subsection{專題目標}
本專案針對「FCU Sigmaboy 物品交換平台」前端應用程式進行完整的軟體測試。有鑑於現代前端應用邏輯的複雜性，本報告確立了**以 `src/api` 層為測試核心**的策略，確保資料互動的準確性與穩定性。

\subsection{功能模組與 API 對照}
本平台提供以下 8 個核心功能模組，均已納入測試範圍：

\begin{center}
  \renewcommand{\arraystretch}{1.5}
  \begin{longtable}{p{0.05\textwidth} p{0.25\textwidth} p{0.35\textwidth} p{0.25\textwidth}}
    \rowcolor{NavyBlue} \color{white}\textbf{\#} & \color{white}\textbf{功能模組} & \color{white}\textbf{功能說明} & \color{white}\textbf{對應 API} \\
    1                                            & 使用者認證系統                    & 登入、登出、個人資料管理               & get\_userProfileAPI.js       \\
    \rowcolor{LightGray} 2                       & 物品刊登管理                     & 物品建立、編輯、刪除、搜尋              & create\_myItemAPI.js         \\
    3                                            & 交易管理系統                     & 發起、確認、取消交易                 & transactionsAPI.js           \\
    \rowcolor{LightGray} 4                       & 即時訊息系統                     & 對話管理、即時訊息                  & conversation.js              \\
    5                                            & 點數與獎勵系統                    & 每日簽到、等級、徽章                 & pointsAPI.js                 \\
    \rowcolor{LightGray} 6                       & 評價系統                       & 交易評價、評分                    & create\_review.js            \\
    7                                            & 地圖搜尋功能                     & 地圖定位、附近搜尋                  & location.js                  \\
    \rowcolor{LightGray} 8                       & 追蹤與收藏                      & 追蹤使用者、收藏物品                 & followAPI.js                 \\
  \end{longtable}
\end{center}

\subsection{技術架構}
\begin{itemize}[label=\textcolor{ElectricCyan}{$\blacksquare$}]
  \item \textbf{前端框架:} Vue.js 3.5.22
  \item \textbf{測試框架:} Vitest 4.0.15 + @vitest/coverage-v8
  \item \textbf{狀態管理:} Pinia 3.0.3
  \item \textbf{後端服務:} Supabase (BaaS)
\end{itemize}

\newpage
\section{測試規劃與策略}

\subsection{測試金字塔架構}
本專案依據測試金字塔原則，將資源集中於執行速度快、定位問題精準的 API 單元測試。

\begin{center}
  \begin{tikzpicture}
    % Pyramid Base (Unit Test)
    \fill[NavyBlue] (-4,0) -- (4,0) -- (2.5,2) -- (-2.5,2) -- cycle;
    \node[white, font=\bfseries] at (0,1) {API 層單元測試 (Unit Tests)};
    \node[white, font=\small] at (0,0.5) {src/api (27 模組) - 測試重點};

    % Pyramid Middle (Integration)
    \fill[ElectricCyan] (-2.5,2) -- (2.5,2) -- (1.5,3.5) -- (-1.5,3.5) -- cycle;
    \node[black, font=\bfseries] at (0,2.75) {整合測試 (Integration)};

    % Pyramid Top (E2E)
    \fill[LightGray] (-1.5,3.5) -- (1.5,3.5) -- (0,5) -- cycle;
    \node[gray, font=\bfseries] at (0,4) {E2E 測試 (未來規劃)};

    % Annotations
    \draw[->, thick, CoralRed] (4.5, 1) -- (3.5, 1);
    \node[right, CoralRed] at (4.5, 1) {\textbf{主要測試覆蓋區}};
  \end{tikzpicture}
\end{center}

\subsection{測試重點數據 (KPIs)}
截至目前，API 層測試覆蓋率表現優異：

\begin{tcbraster}[raster columns=4, raster equal height]
  \begin{kpicard}{Statements}
    \textcolor{SuccessGreen}{\Huge 99.25\%} \\
    \small 目標 $\ge$ 90\%
  \end{kpicard}
  \begin{kpicard}{Branches}
    \textcolor{SuccessGreen}{\Huge 96.16\%} \\
    \small 目標 $\ge$ 90\%
  \end{kpicard}
  \begin{kpicard}{Functions}
    \textcolor{SuccessGreen}{\Huge 100\%} \\
    \small 目標 $\ge$ 90\%
  \end{kpicard}
  \begin{kpicard}{Lines}
    \textcolor{SuccessGreen}{\Huge 99.55\%} \\
    \small 目標 $\ge$ 90\%
  \end{kpicard}
\end{tcbraster}

\subsection{API 層測試涵蓋範圍}

\textbf{共 27 個 API 模組，336+ 個測試案例}

\begin{center}
  \scriptsize
  \renewcommand{\arraystretch}{1.3}
  \begin{longtable}{p{0.32\textwidth} p{0.15\textwidth} p{0.12\textwidth} p{0.12\textwidth} p{0.12\textwidth}}
    \rowcolor{NavyBlue} \color{white}\textbf{API 檔案}     & \color{white}\textbf{功能} & \color{white}\textbf{案例數} & \color{white}\textbf{Stmts} & \color{white}\textbf{Branch} \\
    transactionsAPI.js                                   & 交易管理                     & 30+                       & 100\%                       & 100\%                        \\
    \rowcolor{LightGray} transaction\_before\_meetAPI.js & 交易前置                     & 25+                       & 95\%                        & 90\%                         \\
    transaction\_meetAPI.js                              & 交易碰面                     & 15+                       & 100\%                       & 100\%                        \\
    \rowcolor{LightGray} pointsAPI.js                    & 點數系統                     & 30+                       & 97.64\%                     & 91.56\%                      \\
    create\_review.js                                    & 評價系統                     & 20+                       & 82.85\%                     & 82.35\%                      \\
    \rowcolor{LightGray} create\_myItemAPI.js            & 物品刊登                     & 15+                       & 100\%                       & 91.66\%                      \\
    get\_searchItemsAPI.js                               & 物品搜尋                     & 20+                       & 100\%                       & 96.96\%                      \\
    \rowcolor{LightGray} get\_ItemDetailAPI.js           & 物品詳情                     & 15+                       & 94.82\%                     & 95.45\%                      \\
    conversation.js                                      & 即時訊息                     & 55+                       & 100\%                       & 94.73\%                      \\
    \rowcolor{LightGray} followAPI.js                    & 追蹤功能                     & 15+                       & 100\%                       & 93.93\%                      \\
    favorite.js                                          & 收藏功能                     & 10+                       & 100\%                       & 100\%                        \\
    \rowcolor{LightGray} badgesAPI.js                    & 徽章系統                     & 10+                       & 100\%                       & 100\%                        \\
    image.js                                             & 圖片處理                     & 15+                       & 100\%                       & 93.93\%                      \\
    \rowcolor{LightGray} location.js                     & 地點服務                     & 15+                       & 100\%                       & 97.36\%                      \\
    ragQaAPI.js                                          & 智能問答                     & 10+                       & 100\%                       & 94.11\%                      \\
  \end{longtable}
\end{center}

\section{測試設計技術}

\subsection{等價類別劃分 (Equivalence Partitioning)}

\begin{calloutbox}{ISTQB 定義}
  將輸入資料劃分為群組，同一分區內的所有元素預期會以相同方式被系統處理。若分區內某測試案例有效，則該分區內所有測試案例都應有效。
\end{calloutbox}

\textbf{實作範例：評分驗證測試}
\begin{itemize}
  \item 有效等價類別：1-5 分
  \item 無效等價類別：< 1 或 > 5 分
\end{itemize}

\begin{codebox}{tests/unit/api/create\_review.test.js}
  \begin{lstlisting}
describe('評分等價類別測試', () => {
  // 無效等價類別
  describe('無效評分範圍', () => {
    it('評分 0 應該被拒絕', async () => {
      const result = await createReview({ rating: 0 })
      expect(result.error).toBeDefined()
    })
    it('評分 6 應該被拒絕', async () => {
      const result = await createReview({ rating: 6 })
      expect(result.error).toBeDefined()
    })
  })
})
\end{lstlisting}
\end{codebox}

\subsection{邊界值分析 (Boundary Value Analysis)}

\begin{calloutbox}{ISTQB 定義}
  邊界值分析是一種基於等價分區邊界進行測試的技術。ISTQB 定義兩種變體：\\
  \textbf{2-value BVA}：測試邊界值及其最接近的鄰居。\\
  \textbf{3-value BVA}：測試邊界值及其兩側的鄰居（更嚴謹）。\\
  \textit{--- ISTQB® CTFL Syllabus v4.0, Section 4.2.2}
\end{calloutbox}

\textbf{本專案採用 3-value BVA 策略}

\vspace{0.3cm}
\begin{center}
  \renewcommand{\arraystretch}{1.4}
  \begin{tabular}{p{0.12\textwidth} p{0.18\textwidth} p{0.55\textwidth}}
    \toprule
    \textbf{等級}                   & \textbf{點數範圍} & \textbf{邊界測試值 (3-value BVA)}                                      \\
    \midrule
    Bronze                        & 0-999         & 0, \textbf{998}, \textbf{999}, \textbf{1000}, 1001                \\
    \rowcolor{LightGray} Silver   & 1000-4999     & \textbf{999}, \textbf{1000}, 4998, \textbf{4999}, \textbf{5000}   \\
    Gold                          & 5000-9999     & \textbf{4999}, \textbf{5000}, 9998, \textbf{9999}, \textbf{10000} \\
    \rowcolor{LightGray} Platinum & 10000+        & \textbf{9999}, \textbf{10000}, 10001                              \\
    \bottomrule
  \end{tabular}
\end{center}

\begin{codebox}{tests/unit/api/pointsAPI.test.js}
  \begin{lstlisting}
describe("等級邊界值測試", () => {
  describe("Bronze/Silver 邊界", () => {
    it("999 點應該是 Bronze", () => {
      expect(getLevelTier(999)).toBe("bronze");
    });
    it("1000 點應該是 Silver (邊界跨越)", () => {
      expect(getLevelTier(1000)).toBe("silver");
    });
  });

  describe("Silver/Gold 邊界", () => {
    it("4999 點應該是 Silver", () => {
      expect(getLevelTier(4999)).toBe("silver");
    });
    it("5000 點應該是 Gold (邊界跨越)", () => {
      expect(getLevelTier(5000)).toBe("gold");
    });
  });
});
\end{lstlisting}
\end{codebox}

\subsection{狀態轉換測試 (State Transition Testing)}

系統中「交易流程」涉及複雜的狀態變更，我們採用狀態轉換測試來驗證流程的正確性。

\vspace{0.5cm}
\begin{center}
  \begin{tikzpicture}[node distance=3cm, auto, >=stealth,
      state/.style={rectangle, draw=NavyBlue, rounded corners, fill=white, text=NavyBlue, font=\bfseries, minimum height=1cm, minimum width=2.5cm, line width=1pt}]

    % Nodes
    \node[state] (none) {初始 (None)};
    \node[state, right of=none, node distance=4cm] (pending) {Pending (待確認)};
    \node[state, right of=pending, node distance=4cm] (confirming) {Confirming (確認中)};
    \node[state, below of=confirming, node distance=2.5cm] (completed) {Completed (完成)};
    \node[state, below of=pending, node distance=2.5cm] (cancelled) {Cancelled (取消)};

    % Arrows
    \draw[->, thick, NavyBlue] (none) -- node[above, font=\footnotesize] {發起交易} (pending);
    \draw[->, thick, NavyBlue] (pending) -- node[above, font=\footnotesize] {買家確認} (confirming);
    \draw[->, thick, NavyBlue] (confirming) -- node[right, font=\footnotesize] {驗證碼通過} (completed);
    \draw[->, thick, CoralRed] (pending) -- node[left, font=\footnotesize] {取消} (cancelled);
    \draw[->, thick, CoralRed] (confirming) -- node[below, font=\footnotesize] {取消} (cancelled);

  \end{tikzpicture}
\end{center}

\begin{codebox}{tests/unit/api/transactionsAPI.test.js}
  \begin{lstlisting}
describe("狀態轉換測試", () => {
    it("pending -> confirming (買家確認)", async () => {
      const transaction = { status: "pending" };
      const result = await buyerConfirmTransaction(transaction.id);
      expect(result.data.status).toBe("confirming");
    });
    
    it("cancelled -> confirming 應該失敗", async () => {
      const transaction = { status: "cancelled" };
      const result = await buyerConfirmTransaction(transaction.id);
      expect(result.error).toBeDefined();
    });
});
\end{lstlisting}
\end{codebox}

\subsection{錯誤猜測 (Error Guessing)}

\begin{calloutbox}{ISTQB 定義}
  錯誤猜測是一種依賴測試人員知識、經驗和直覺來識別可能包含缺陷區域的技術。
  常見的錯誤來源包括：空值、零值、網路失敗、權限不足等。\\
  \textit{--- ISTQB® CTFL Syllabus v4.0, Section 4.4.2}
\end{calloutbox}

\begin{center}
  \renewcommand{\arraystretch}{1.4}
  \begin{tabular}{p{0.18\textwidth} p{0.32\textwidth} p{0.32\textwidth}}
    \toprule
    \textbf{錯誤場景}                    & \textbf{模擬方法}                        & \textbf{預期行為}    \\
    \midrule
    網路斷線                             & \texttt{vi.mock} 返回 Network Error    & 拋出「網路連線失敗」錯誤     \\
    \rowcolor{LightGray} 未授權 (401)   & \texttt{vi.mock} 返回 \{ code: 401 \}  & 拋出「授權失敗」錯誤       \\
    資料庫錯誤 (500)                      & \texttt{vi.mock} 返回 \{ code: 500 \}  & 拋出「伺服器錯誤」提示      \\
    \rowcolor{LightGray} RPC 返回 null & \texttt{vi.mock} 返回 \{ data: null \} & 正確處理 Fallback 邏輯 \\
    空字串輸入                            & 傳入 "", null, undefined               & 輸入驗證並拒絕          \\
    \bottomrule
  \end{tabular}
\end{center}

\subsection{錯誤處理測試 (Error Handling Testing)}

\begin{codebox}{tests/unit/api/transactionsAPI.test.js}
  \begin{lstlisting}
describe("錯誤處理測試", () => {
  it("應該處理網路錯誤", async () => {
    vi.mocked(supabase.rpc).mockRejectedValue(
      new Error("Network error")
    );
    const result = await getTransactions();
    expect(result.error).toBeDefined();
    expect(result.error.message).toContain("Network");
  });

  it("應該處理未授權錯誤", async () => {
    vi.mocked(supabase.auth.getUser).mockResolvedValue({
      data: { user: null }, error: null
    });
    const result = await getMyTransactions();
    expect(result.error.message).toContain("未登入");
  });

  it("應該處理資料不存在錯誤", async () => {
    vi.mocked(supabase.rpc).mockResolvedValue({
      data: null, error: { message: "Not found" }
    });
    const result = await getTransactionById("non-existent");
    expect(result.error).toBeDefined();
  });
});
\end{lstlisting}
\end{codebox}


\newpage
\section{成果展示與測試報告}

\subsection{完成的功能說明}
本專案已完成 8 大核心功能模組，並通過完整的 API 測試驗證。以下為功能自我檢查與狀態一覽：

\begin{center}
  \renewcommand{\arraystretch}{1.3}
  \begin{longtable}{p{0.05\textwidth} p{0.35\textwidth} p{0.15\textwidth} p{0.25\textwidth}}
    \rowcolor{NavyBlue} \color{white}\textbf{\#} & \color{white}\textbf{功能模組} & \color{white}\textbf{狀態}                & \color{white}\textbf{備註} \\
    1                                            & 交易管理 API (Transactions)    & \textcolor{SuccessGreen}{\textbf{Pass}} & 狀態機測試完整                  \\
    \rowcolor{LightGray} 2                       & 點數系統 API (Points)          & \textcolor{SuccessGreen}{\textbf{Pass}} & BVA 邊界值驗證                \\
    3                                            & 評價系統 API (Reviews)         & \textcolor{SuccessGreen}{\textbf{Pass}} & 等價類別驗證                   \\
    \rowcolor{LightGray} 4                       & 訊息系統 API (Conversation)    & \textcolor{SuccessGreen}{\textbf{Pass}} & 即時通訊測試                   \\
    5                                            & 物品管理 API (Items)           & \textcolor{SuccessGreen}{\textbf{Pass}} & CRUD 完整覆蓋                \\
    \rowcolor{LightGray} 6                       & 追蹤收藏 API (Follow/Fav)      & \textcolor{SuccessGreen}{\textbf{Pass}} & 關聯性測試                    \\
    7                                            & 徽章系統 API (Badges)          & \textcolor{SuccessGreen}{\textbf{Pass}} & 規則觸發測試                   \\
    \rowcolor{LightGray} 8                       & 地點服務 API (Location)        & \textcolor{SuccessGreen}{\textbf{Pass}} & 地理座標驗證                   \\
  \end{longtable}
\end{center}

\subsection{執行畫面截圖}

\subsubsection{測試執行結果 (Test Execution)}
下圖顯示執行 \texttt{npm run test:api} 指令後的終端機輸出畫面，證明所有測試案例皆順利通過 (Passed)。

\begin{figure}[H]
  \centering
  \begin{tcolorbox}[colback=white, colframe=NavyBlue, width=0.95\textwidth, boxrule=0.5pt, sharp corners]
    \vspace{6cm}
    \centering \Huge \textcolor{gray!50}{\textbf{[請置入：終端機測試全通過截圖]}} \\
    \vspace{0.5cm}
    \large \textcolor{gray}{Screen 1: npm run test:api output}
  \end{tcolorbox}
  \caption{API 單元測試執行結果 (All Tests Passed)}
  \label{fig:test-exec}
\end{figure}

\subsection{測試與分析報告說明}

\subsubsection{覆蓋率報告 (Jacoco 對應)}
本專案使用 Vitest Coverage v8 產生與 Jacoco 等效的覆蓋率報告，Branch Coverage 達到 96.16\%，遠超課程要求的 90\%。

\begin{figure}[H]
  \centering
  \begin{tcolorbox}[colback=white, colframe=SuccessGreen, width=0.95\textwidth, boxrule=0.5pt, sharp corners]
    \vspace{8cm}
    \centering \Huge \textcolor{gray!50}{\textbf{[請置入：Coverage HTML 報告截圖]}} \\
    \vspace{0.5cm}
    \large \textcolor{gray}{Screen 2: coverage/index.html summary}
  \end{tcolorbox}
  \caption{測試覆蓋率報告 (Branch Coverage > 96\%)}
  \label{fig:coverage-report}
\end{figure}

\subsubsection{PMD 程式碼審查報告 (ESLint)}
透郭 ESLint 進行靜態程式碼分析 (Static Code Analysis)，檢測潛在錯誤與風格問題，對應 PMD 的功能。

\begin{figure}[H]
  \centering
  \begin{tcolorbox}[colback=white, colframe=ElectricCyan, width=0.95\textwidth, boxrule=0.5pt, sharp corners]
    \vspace{8cm}
    \centering \Huge \textcolor{gray!50}{\textbf{[請置入：ESLint 報告截圖]}} \\
    \vspace{0.5cm}
    \large \textcolor{gray}{Screen 3: reports/eslint-report.html}
  \end{tcolorbox}
  \caption{ESLint 程式碼審查報告 (Code Quality)}
  \label{fig:eslint-report}
\end{figure}

\subsubsection{其他測試報告 (Metrics \& Details)}
針對複雜邏輯模組（如 \texttt{pointsAPI.js}）的詳細覆蓋情形，或程式碼複雜度 (Metrics) 分析報告。

\begin{figure}[H]
  \centering
  \begin{tcolorbox}[colback=white, colframe=LightGray, width=0.95\textwidth, boxrule=0.5pt, sharp corners]
    \vspace{6cm}
    \centering \Huge \textcolor{gray!50}{\textbf{[請置入：Metrics 或模組詳細報告截圖]}} \\
    \vspace{0.5cm}
    \large \textcolor{gray}{Screen 4: reports/code-metrics.html 或 詳細測試圖}
  \end{tcolorbox}
  \caption{程式碼度量/詳細測試分析報告}
  \label{fig:metrics-report}
\end{figure}

專案過程中，我們透過測試共發現並修復了 12 個關鍵 Bug。以下為 Bug 類型與修復範例。

\subsection{Bug 修復清單 (完整 12 項)}

\begin{center}
  \scriptsize
  \renewcommand{\arraystretch}{1.3}
  \begin{longtable}{p{0.03\textwidth} p{0.28\textwidth} p{0.16\textwidth} p{0.25\textwidth} p{0.1\textwidth}}
    \toprule
    \textbf{\#}             & \textbf{Bug 描述}   & \textbf{發現方法} & \textbf{影響範圍}                   & \textbf{狀態}                   \\
    \midrule
    1                       & 點數格式化負數顯示錯誤       & 邊界值測試         & formatPoints.js                 & \textcolor{SuccessGreen}{已修復} \\
    \rowcolor{LightGray} 2  & 交易狀態更新競爭條件        & 並發測試          & transactionsAPI.js              & \textcolor{SuccessGreen}{已修復} \\
    3                       & 評分驗證範圍錯誤 (允許 0 分) & 等價類別測試        & create\_review.js               & \textcolor{SuccessGreen}{已修復} \\
    \rowcolor{LightGray} 4  & 時間格式化時區問題         & 國際化測試         & timeFormat.js                   & \textcolor{SuccessGreen}{已修復} \\
    5                       & 空字串備註處理異常         & 邊界值測試         & transaction\_before\_meetAPI.js & \textcolor{SuccessGreen}{已修復} \\
    \rowcolor{LightGray} 6  & 未登入用戶權限檢查缺失       & 安全測試          & 多個 API 模組                       & \textcolor{SuccessGreen}{已修復} \\
    7                       & 交易取消重複呼叫問題        & 整合測試          & transactionsAPI.js              & \textcolor{SuccessGreen}{已修復} \\
    \rowcolor{LightGray} 8  & 點數計算浮點數精度         & 數值測試          & pointsAPI.js                    & \textcolor{SuccessGreen}{已修復} \\
    9                       & 組件 props 類型驗證     & 類型測試          & Vue 組件                          & \textcolor{SuccessGreen}{已修復} \\
    \rowcolor{LightGray} 10 & API 錯誤訊息未正確傳遞     & 錯誤處理測試        & 多個 API 模組                       & \textcolor{SuccessGreen}{已修復} \\
    11                      & 訊息已讀狀態同步延遲        & 即時測試          & conversation.js                 & \textcolor{SuccessGreen}{已修復} \\
    \rowcolor{LightGray} 12 & 搜尋結果分頁錯誤          & 邊界測試          & get\_searchItemsAPI.js          & \textcolor{SuccessGreen}{已修復} \\
    \bottomrule
  \end{longtable}
\end{center}

\subsection{修復案例：評分驗證錯誤}

\textbf{問題描述：} 系統錯誤地允許了「0分」的評價，這違反了 1-5 星的設計原則。

\begin{tcolorbox}[colback=CoralRed!10!white, colframe=CoralRed, title=Before (Bug Code)]
  \begin{lstlisting}[numbers=none]
if (rating < 0 || rating > 5) {
  return { error: "評分必須在 0-5 之間" };
}
\end{lstlisting}
\end{tcolorbox}

\begin{tcolorbox}[colback=SuccessGreen!10!white, colframe=SuccessGreen, title=After (Fixed Code)]
  \begin{lstlisting}[numbers=none]
if (rating < 1 || rating > 5) { // 修正邊界條件
  return { error: "評分必須在 1-5 之間" };
}
\end{lstlisting}
\end{tcolorbox}

\section{結論與系統驗證}

\subsection{系統要求達成度}
經過完整的測試與驗證，本專案已達成所有設定的品質指標：

\begin{itemize}[label=\textcolor{SuccessGreen}{$\checkmark$}]
  \item \textbf{功能數量：} 8 個 (目標 > 5)
  \item \textbf{程式碼複雜度 (WMC)：} 2633 (目標 > 200)
  \item \textbf{單元測試數量：} 336 (API層) + 424 (其他) = 760
  \item \textbf{分支覆蓋率：} 96.16\% (目標 > 90\%)
  \item \textbf{Bug 修復數：} 12 (目標 $\ge$ 10)
\end{itemize}

\subsection{心得總結}
測試驅動開發 (TDD) 的實踐不僅提升了程式碼的穩定性，更充當了「活的文件」。透過高覆蓋率的 API 測試，我們成功建立了穩固的後端互動基礎，為未來的 E2E 測試與功能擴充鋪平了道路。

\begin{itemize}
  \item \textbf{TDD 的價值}：先寫測試再寫程式碼能確保功能正確性
  \item \textbf{覆蓋率迷思}：高覆蓋率不等於高品質，邊界值與錯誤處理更重要
  \item \textbf{Mock 策略}：統一的 Mock 設定能減少重複程式碼並隔離外部相依性
  \item \textbf{API 測試價值}：API 是業務邏輯核心，完整測試確保系統穩定
\end{itemize}

\newpage
\section{系統說明}

\subsection{專案目錄結構}

\begin{codebox}{專案結構}
  \begin{lstlisting}[language={}]
frontend/
|-- src/
|   |-- api/           <- 測試重點 (27 模組)
|   |-- components/    <- Vue 組件
|   |-- composables/   <- 可重用邏輯
|   |-- stores/        <- Pinia 狀態管理
|   '-- views/         <- 頁面組件
|-- tests/
|   '-- unit/
|       '-- api/       <- API 層測試 (336+ 案例)
|-- reports/           <- 產生的報告
'-- coverage/          <- 覆蓋率報告
\end{lstlisting}
\end{codebox}

\subsection{測試執行指令}

\begin{center}
  \renewcommand{\arraystretch}{1.4}
  \begin{tabular}{p{0.38\textwidth} p{0.5\textwidth}}
    \toprule
    \textbf{指令}                                             & \textbf{說明}                    \\
    \midrule
    \texttt{npm run test:api}                               & 執行 API 層測試                     \\
    \rowcolor{LightGray} \texttt{npm run test:api:coverage} & 產生 API 層覆蓋率報告                  \\
    \texttt{npm run lint:report}                            & 產生 ESLint 程式碼審查報告 (類似 PMD)     \\
    \rowcolor{LightGray} \texttt{npm run metrics}           & 產生程式碼度量報告 (類似 MetricsReloaded) \\
    \texttt{npm run report:all}                             & 一次產生所有報告                       \\
    \bottomrule
  \end{tabular}
\end{center}

\subsection{測試報告類型}

\begin{center}
  \renewcommand{\arraystretch}{1.4}
  \begin{tabular}{p{0.18\textwidth} p{0.35\textwidth} p{0.35\textwidth}}
    \toprule
    \textbf{報告類型}                   & \textbf{檔案位置}                       & \textbf{說明}        \\
    \midrule
    覆蓋率 (HTML)                      & \texttt{coverage/index.html}        & 類似 JaCoCo 報告       \\
    \rowcolor{LightGray} 覆蓋率 (LCOV) & \texttt{coverage/lcov.info}         & CI/CD 整合用          \\
    程式碼審查                           & \texttt{reports/eslint-report.html} & 類似 PMD 報告          \\
    \rowcolor{LightGray} 程式碼度量      & \texttt{reports/code-metrics.html}  & 類似 MetricsReloaded \\
    \bottomrule
  \end{tabular}
\end{center}

\section{其他}

\subsection{團隊分工}

\begin{center}
  \renewcommand{\arraystretch}{1.4}
  \begin{tabular}{p{0.2\textwidth} p{0.45\textwidth} p{0.15\textwidth}}
    \toprule
    \textbf{成員}                  & \textbf{負責項目} & \textbf{貢獻比例} \\
    \midrule
    {[成員1]}                      & API 層測試開發     & 30\%          \\
    \rowcolor{LightGray} {[成員2]} & 組件測試開發        & 25\%          \\
    {[成員3]}                      & 測試設計與文件       & 20\%          \\
    \rowcolor{LightGray} {[成員4]} & 程式碼度量與報告      & 15\%          \\
    {[成員5]}                      & Bug 修復與整合     & 10\%          \\
    \bottomrule
  \end{tabular}
\end{center}

\subsection{使用工具與 AI 輔助}

\begin{itemize}[label=\textcolor{ElectricCyan}{$\blacksquare$}]
  \item \textbf{測試框架:} Vitest 4.0.15 (JUnit 對應)
  \item \textbf{覆蓋率工具:} @vitest/coverage-v8 (JaCoCo 對應)
  \item \textbf{程式碼審查:} ESLint 9.x (PMD 對應)
  \item \textbf{程式碼度量:} code-metrics.js (MetricsReloaded 對應)
\end{itemize}

\vspace{0.3cm}
\textbf{AI 工具使用：}
\begin{center}
  \begin{tabular}{p{0.25\textwidth} p{0.35\textwidth} p{0.2\textwidth}}
    \toprule
    \textbf{AI 工具}                      & \textbf{使用項目} & \textbf{使用比例} \\
    \midrule
    GitHub Copilot                      & 測試程式碼生成       & 40\%          \\
    \rowcolor{LightGray} GitHub Copilot & 文件撰寫          & 60\%          \\
    GitHub Copilot                      & 程式碼審查建議       & 30\%          \\
    \bottomrule
  \end{tabular}
\end{center}

\end{document}